\documentclass[11pt]{article}

\usepackage[T1]{fontenc}
\usepackage{lmodern}
\usepackage{amsmath,amssymb,amsthm}
\usepackage[margin=1in]{geometry}
\usepackage{enumitem}
\usepackage{xcolor}
\usepackage{microtype}
\usepackage[hidelinks]{hyperref}

\newcommand{\R}{\mathbb{R}}
\newcommand{\pos}[1]{\left(#1\right)_{+}}

\title{Paper 2, Theorem 1.3(iv): Exponent-Domain Amendment}
\author{Internal research note for the \texttt{ShenWork} formalization}
\date{13 July 2026}

\begin{document}
\maketitle

\begin{abstract}
The proof of Theorem 1.3(iv) in Chen--Ruau--Shen invokes the weighted elliptic
estimate of Proposition 2.2 at an exponent which is not necessarily above
one.  The missing admissibility condition is
\(q_* > 2-2m\), where \(q_*=\max\{1,N\alpha/2\}\).  This note derives the
condition, gives an explicit admitted parameter wedge where the printed
argument fails, and records the corrected theorem proved in Lean.  The issue
is a proof gap; no counterexample to the boundedness conclusion is claimed.
\end{abstract}

\section{Source and printed alternative}

The source examined is Le Chen, Ian Ruau, and Wenxian Shen,
\emph{Chemotaxis models with signal-dependent sensitivity and a logistic-type
source, I: Boundedness and global existence}, arXiv:2512.14858v1, 16 December
2025.  The local archival copy is \path{paper2.pdf}.

Theorem 1.3(iv) assumes
\begin{equation}
  \beta\geq\frac12,
  \qquad
  \alpha=2m+\gamma-2,
  \label{eq:critical}
\end{equation}
together with the smallness alternative (1.25).  It is stated for all
positive \(m,\gamma,\alpha\).

\section{The missing exponent-domain condition}

Section 5.4 invokes Proposition 2.2 at
\begin{equation}
  s(P)=\frac{P+\alpha}{\gamma}.
  \label{eq:sP}
\end{equation}
Proposition 2.2 requires \(s(P)>1\).  Substituting
\eqref{eq:critical} gives
\begin{align*}
  \frac{P+\alpha}{\gamma}>1
  &\iff P+2m+\gamma-2>\gamma\\
  &\iff P>2-2m.
\end{align*}
The proof selects a seed exponent immediately to the right of
\begin{equation}
  q_*:=\max\left\{1,\frac{N\alpha}{2}\right\}.
\end{equation}
It therefore needs
\begin{equation}
  \boxed{q_*>2-2m,}
  \label{eq:missing}
\end{equation}
or equivalently \((q_*+\alpha)/\gamma>1\).  Condition
\eqref{eq:missing} is not implied by the printed hypotheses.  The claim in
Section 5.4 that the exponent supplied to Proposition 2.2 is above one for
all seed powers is false in the remaining wedge.

\section{Concrete admitted parameter wedge}

On the interval choose
\begin{equation}
  N=1,\qquad m=\frac14,\qquad \gamma=\frac72,\qquad
  \alpha=2,\qquad \beta=1.
\end{equation}
Then \(\alpha=2m+\gamma-2\), while
\begin{equation}
  q_*=\max\{1,\alpha/2\}=1,
  \qquad
  s(q_*)=\frac{1+2}{7/2}=\frac67<1.
\end{equation}
Furthermore \(\pos{N\alpha-2}=0\).  The first disjunct in (1.25) is
therefore automatic, so the printed theorem places no smallness condition on
\(\chi_0\).  Proposition 2.2 becomes available only for \(P>3/2\).  At such
powers the squared-chemotaxis coefficient is nonzero, and the automatic first
disjunct of (1.25) provides no estimate producing the required low seed.

This demonstrates a gap in the written proof, not a counterexample to the
boundedness statement.  A new estimate may recover part or all of the
uncovered parameters.

\section{Proof-supported amendment}

The Section 5.4 route is valid after adding \eqref{eq:missing}.  On the
one-dimensional interval this reads
\begin{equation}
  \max\{1,\alpha/2\}>2-2m.
\end{equation}
The literal profile defining (1.19) is then continuous at \(q_*\).  The
smallness condition (1.25) yields \(p_0>q_*\) with strict damping.  Since
\(p_0>\alpha/2\), a one-dimensional fixed-target interpolation absorbs
\(\int u^{P+\alpha}\) at one finite target \(P\) above both \(m\) and
\(\gamma\).  A restarted Neumann-semigroup estimate then gives the supremum
bound.  No infinite Moser iteration is needed.

\medskip
\noindent
\fbox{%
\begin{minipage}{0.92\textwidth}
\textbf{Recommended amendment to alternative (iv).}
In addition to the printed hypotheses, assume
\(\max\{1,N\alpha/2\}>2-2m\).  Then the proof in Section 5.4 is valid.
The original statement outside this range remains unproved by the written
argument.
\end{minipage}}

\section{Lean realization}

The formalization keeps the missing condition explicit.
\begin{itemize}[leftmargin=2em]
\item \path{IntervalDomainTheorem13CriticalConstants.lean} realizes the
literal \(M^*\), \(q_*\), and right-sided lower limit \(K\).
\item \path{IntervalDomainTheorem13CriticalSeed.lean} proves the case-(iv)
energy estimate with \(2-2m<P\).
\item \path{IntervalDomainTheorem13CriticalThreshold.lean} obtains the low
seed from (1.25) under \(q_*>2-2m\).
\item \path{IntervalDomainTheorem13CriticalBootstrap.lean} proves
\texttt{boundedBefore\_critical\_case\_iv\_corrected} on the faithful
general-\(m\) interval model by a single fixed-target interpolation and the
proved restarted endpoint.
\end{itemize}
All exported capstones are \texttt{sorry}-free and depend only on
\texttt{propext}, \texttt{Classical.choice}, and \texttt{Quot.sound}.

Alternative (iii) does not share this defect: under
\(\alpha=m+\gamma-1\), the inequality
\((q_*+\alpha)/\gamma>1\) follows from \(q_*\geq1\) and \(m>0\).

\section*{Reference}

L. Chen, I. Ruau, and W. Shen, \emph{Chemotaxis models with signal-dependent
sensitivity and a logistic-type source, I: Boundedness and global existence},
arXiv:2512.14858v1 (2025),
\url{https://arxiv.org/abs/2512.14858}.

\end{document}
